\section{Design Evidence}
\label{sec:motivation}

CoMem does not require a claim that a model ``understands'' a document at one
particular layer. We use three empirical observations to motivate the interface:
information can be accessible before native final-head readout; resuming too
deep creates a sharp fidelity loss; and the dominant bottleneck can move between
retrieval and residual readout depending on the task.

\subsection{Accessibility precedes native readout}
Under our probe protocols, semantic labels become linearly accessible before a
fixed final norm and LM head can express the corresponding native answer
(Figure~\ref{fig:depth}). This ordering appears across the tested Qwen, Llama,
and OLMo backbones. It is readout-dependent evidence, not a causal boundary or
a statement that later layers are unnecessary.

\begin{figure}[t]
\centering
\includegraphics[width=\columnwidth]{figures/depth_motiv.pdf}
\caption{\textbf{Intermediate accessibility motivates, but does not validate,
depth reuse.} (a) Zero-shot task readout remains usable only up to a
depth-dependent boundary. (b) Frozen linear probes become informative earlier
than native LM-head readout under the specified protocols. Exact definitions and
values are in Appendix~\ref{app:ablations}.}
\label{fig:depth}
\end{figure}

\subsection{The split depth is a fidelity knob}
A frozen residual interface degrades as the split moves deeper. Separate
self-distilled adapters shift the usable frontier but do not eliminate it:
RULER changes from 98.29 at $j{=}6$ to 96.07 at $j{=}12$ and 55.41 at
$j{=}18$, while fixed-pack Read speedup rises from $1.166\times$ to
$1.403\times$ and $1.807\times$ (Appendix
Table~\ref{tab:distilled-depth-curve}). Because adapter spans and parameter
counts differ, this is a deployment frontier, not a compute-matched causal
depth ablation.

\subsection{Retrieval and readout are separate failure modes}
On a lexical single-needle diagnostic, BM25 and oracle selection retain the
evidence while recency and reader-attention heuristics deteriorate with length
(Figure~\ref{fig:read}). That observation is deliberately narrow. On LongEval,
BM25 recall is 97--100\% across 8k--128k and matched $j{=}0$ replay is nearly
perfect conditional on a retrieval hit, whereas frozen $j{=}12$ fails and the
distilled interface recovers only 64--76\% hit-conditional accuracy. On
BABILong, both retrieval misses and residual readout contribute, with the
balance varying by task. Thus a bounded memory needs both a suitable selector
and a compatible stored representation.

\begin{figure}[t]
\centering
\includegraphics[width=\columnwidth]{figures/read_motiv.pdf}
\caption{\textbf{Selection is one bottleneck, not always the bottleneck.}
Oracle and BM25 remain strong on the single-needle diagnostic; other analyses
separate evidence recall from reader accuracy and show substantial residual
readout loss on multi-step tasks.}
\label{fig:read}
\end{figure}

\subsection{Implications}
The evidence suggests three design requirements. First, expose split depth as a
measured quality--cost knob. Second, retrieve a bounded relevant pack but allow
full upper-layer interaction among its chunks. Third, preserve enough
lower-layer document context during Write that the stored residual is compatible
with later composition. Section~\ref{sec:method} instantiates the first two;
Section~\ref{sec:experiments} evaluates independent and overlapping Write
variants for the third.
